\DocumentMetadata{
  pdfversion=2.0,
  pdfstandard=ua-2,
  lang=en-US,
  tagging=on,
  tagging-setup={math/setup=mathml-SE,tabsorder=structure}
}
\documentclass[11pt,letterpaper]{article}
\usepackage[margin=0.68in,includefoot]{geometry}
\usepackage{newcomputermodern}
\usepackage{amsmath,amscd,mathtools}
\usepackage{tikz,adjustbox}
\providecommand{\square}{\mdlgwhtsquare}
\usepackage[hidelinks]{hyperref}
\hypersetup{
  pdftitle={Analysis PhD Exam, August 2018},
  pdfauthor={University of Florida Department of Mathematics},
  pdfsubject={Graduate mathematics examination},
  pdfdisplaydoctitle=true,
  bookmarksopen=true
}
\setcounter{secnumdepth}{0}
\setlength{\parindent}{0pt}
\setlength{\parskip}{0.28em}
\setlength{\emergencystretch}{3em}
\setlength{\leftmargini}{2.15em}
\setlength{\leftmarginii}{2.65em}
\setlength{\leftmarginiii}{3em}
\pagestyle{empty}
\raggedbottom
\allowdisplaybreaks
\NewTaggingSocketPlug{block/list/label}{examproblem}{%
  \tagstructbegin{tag=Lbl}\tagstructend%
  \tagstructbegin{tag=\LBody}%
  \tagstructbegin{tag=\ProblemHeadingTag,title-o={\ProblemHeadingText}}%
  \tagmcbegin{tag=\ProblemHeadingTag,actualtext={\ProblemHeadingText}}%
  #2%
  \tagmcend%
  \tagstructend%
}
\newenvironment{examproblems}[2]{%
  \def\ProblemHeadingTag{#2}%
  \AssignTaggingSocketPlug{block/list/label}{examproblem}%
  \begin{enumerate}%
  \setcounter{enumi}{#1}%
  \renewcommand{\labelenumi}{\textbf{\theenumi.}}%
  \setlength{\topsep}{0.35em}%
  \setlength{\partopsep}{0pt}%
  \setlength{\itemsep}{0.48em}%
  \setlength{\parsep}{0.18em}%
}{\end{enumerate}}
\newenvironment{examsubparts}{%
  \AssignTaggingSocketPlug{block/list/label}{default}%
  \begin{enumerate}%
  \setlength{\topsep}{0.18em}%
  \setlength{\partopsep}{0pt}%
  \setlength{\itemsep}{0.16em}%
  \setlength{\parsep}{0.1em}%
}{\end{enumerate}}
\newenvironment{examinstructions}{%
  \par\begingroup\itshape
}{%
  \par\endgroup\vspace{0.18em}
}
\makeatletter
\renewcommand\subsection{\@startsection{subsection}{2}{0pt}%
  {-1.25ex plus -.25ex minus -.1ex}%
  {0.55ex plus .1ex}%
  {\normalfont\large\bfseries}}
\renewcommand{\@maketitle}{%
  \newpage\null\vskip -1em%
  {\footnotesize\bfseries
    \href{https://www.ufl.edu/}{University of Florida}, \href{https://math.ufl.edu/}{Department of Mathematics}\par}%
  \vskip -0.5em%
  \tagstructbegin{tag=H1,title={Analysis PhD Exam, August 2018}}%
  \tagpdfparaOff%
  \tagmcbegin{tag=H1}%
    {\LARGE\bfseries\href{https://gma.math.ufl.edu/exams/analysis-phd/}{Analysis PhD Exam}, August 2018\par}%
  \tagmcend%
  \tagpdfparaOn%
  \tagstructend%
  \vskip 0.25em%
  \tagmcbegin{artifact}%
    \hrule height 1pt%
  \tagmcend%
  \vskip 1em%
}
\makeatother
\title{Analysis PhD Exam, August 2018}
\author{}
\date{}
\begin{document}
\maketitle
\thispagestyle{empty}
\begin{examinstructions}
Do six of eight. Write solutions in a neat and logical fashion, giving complete reasons for all steps.
\end{examinstructions}
\begin{examproblems}{0}{H2}
\def\ProblemHeadingText{Problem 1.}
\item
Let \(X,Y\) be topological spaces. Prove, if \(f:X\to Y\) is continuous, then~\(f\) is Borel measurable.
\def\ProblemHeadingText{Problem 2.}
\item
Let \((X,\mathcal M,\mu)\) be a~\(\sigma\)-finite measure space and let \(f\in L^1(\mu)\). Prove that for every \(\epsilon>0\) there exists a \(\delta>0\) such that if \(E\in\mathcal M\) and \(\mu(E)<\delta\), then
\[
\int_E |f|\,d\mu<\epsilon.
\]
\def\ProblemHeadingText{Problem 3.}
\item
Do two. In each case, give a brief justification for your answer.
\begin{examsubparts}
\item[\textnormal{(a)}]
Give an example, if possible, of a sequence \((f_n)\subset L^1(\mathbb R)\) and an \(f\in L^1(\mathbb R)\) such that \(f_n\to f\) uniformly on \(\mathbb R\) but \(f_n\) does not converge to~\(f\) in the \(L^1\) norm.
\item[\textnormal{(b)}]
Evaluate
\[
\lim_{n\to\infty}\int_1^\infty \frac{e^{inx}}{x^n}\,dx.
\]
\item[\textnormal{(c)}]
Give an example to show that the~\(\sigma\)-finiteness hypothesis is necessary in Tonelli’s theorem.
\end{examsubparts}
\def\ProblemHeadingText{Problem 4.}
\item
Let \((X,\mathcal M,\mu)\) be a~\(\sigma\)-finite measure space and \(\mathcal N\subset\mathcal M\) a sub-\(\sigma\)-algebra. Put \(\nu=\mu|_{\mathcal N}\). Prove, if~\(\nu\) is~\(\sigma\)-finite and \(f\in L^1(\mu)\), then there exists a \(g\in L^1(\nu)\) such that for all \(E\in\mathcal N\),
\[
\int_E f\,d\mu=\int_E g\,d\nu.
\]
\def\ProblemHeadingText{Problem 5.}
\item
Let \((X,\mathcal M,\mu)\) be a~\(\sigma\)-finite measure space. Fix a measurable function \(f:X\to\mathbb C\).
\begin{examsubparts}
\item[\textnormal{(a)}]
Assuming \(f\in L^\infty(\mu)\), show if \(g\in L^1(\mu)\), then \(fg\in L^1(\mu)\) and the mapping \(M_f:L^1(\mu)\to L^1(\mu)\) defined by \(M_fg=fg\) is a bounded linear transformation.
\item[\textnormal{(b)}]
Conversely, show, if \(fg\in L^1(\mu)\) for each \(g\in L^1(\mu)\), then \(f\in L^\infty(\mu)\).
\end{examsubparts}
\def\ProblemHeadingText{Problem 6.}
\item
Fix \(1<p<\infty\) and let \(\frac1p+\frac1q=1\). Let \((f_n)\) be a sequence in \(L^p[0,1]\). Suppose that for every \(g\in L^q[0,1]\), the limit
\[
\lim_{n\to\infty}\int_0^1 f_ng
\]
 exists. Prove that there exists an \(f\in L^p[0,1]\) such that
\[
\lim_{n\to\infty}\int_0^1 f_ng=\int_0^1 fg
\]
 for all \(g\in L^q[0,1]\). (Hint: first show that the assignment \(g\mapsto\lim\int f_ng\) defines a linear functional on \(L^q\).)
\def\ProblemHeadingText{Problem 7.}
\item
Let~\(X\) be a Banach space. Say a sequence \((x_n)\) from~\(X\) converges weakly to \(x\in X\) if \(f(x_n)\to f(x)\) for every \(f\in X^*\). Prove that if~\(Y\) is a closed subspace of~\(X\), \((x_n)\subset Y\), and \(x_n\to x\) weakly, then \(x\in Y\).
\def\ProblemHeadingText{Problem 8.}
\item
Let~\(X\) be a Banach space, \(Y\subset X\) a closed subspace. Say~\(Y\) is complemented in~\(X\) if there exists a closed subspace \(Z\subset X\) such that \(Y\cap Z=\{0\}\) and \(Y+Z=X\).

\par
Prove, if~\(Y\) is complemented in~\(X\), then there exists a bounded linear operator \(T:X\to Y\) such that \(T(y)=y\) for all \(y\in Y\).
\end{examproblems}
\end{document}
